\section*{Exercises}
\begin{exercise}\label{ex:17.1}Explain why infinitely many sequence indices inside a ball do not imply infinitely many distinct points inside that ball. Does the proof need distinct points?\end{exercise}
\begin{exercise}\label{ex:17.2}An assistant chooses an arbitrary $n_k\in S_k$ without requiring $n_k>n_{k-1}$. Identify the gap and repair the choice. Explain why the repair is always possible.\end{exercise}
\begin{exercise}\label{ex:17.3}Prove that the discrete metric on $\NN$ is complete but not totally bounded. Use it to explain the difference between completeness and the conclusion of Theorem~\ref{thm:subsequence}.\end{exercise}
\begin{exercise}\label{ex:17.4}Give a full finite-ball-cover argument for total boundedness of $[a,b]\subset\RR$, where $a<b$. State exactly what the number of balls may depend on.\end{exercise}
\begin{exercise}\label{ex:17.5}A directed graph has edges $v\to u$, $u\to w$, and $v\to w$. Compare the two-edge-reachability and exact-second-neighborhood definitions at $v$. Explain the formalization error caused by replacing one with the other.\end{exercise}
\begin{exercise}\label{ex:17.6}Write a minimal-background request for the contraction theorem and one for the subsequence theorem. For each, identify a concept needed in that proof but not in the other. Then state what your repository reading record may legitimately claim without a completed repository build.\end{exercise}
